\documentclass[10pt,a4paper]{amsart}
\usepackage[margin=19mm]{geometry}
\usepackage{amsmath,amssymb,mathtools}
\usepackage[scr=rsfs,cal=euler]{mathalfa}
\usepackage{xfrac}
\usepackage[unicode,hidelinks,bookmarks=false]{hyperref}
\newcommand{\ie}{i.e.\ }
\newcommand{\cf}{cf.\ }
\newcommand{\resp}{respectively\ }
\newcommand{\Subsection}{\subsection}
\providecommand{\nobreakdash}{\nobreak\hbox{-}}
\setlength{\emergencystretch}{2em}
\multlinegap0pt
\usepackage{mathtools}
\usepackage{tensor}
\usepackage{amssymb}
\usepackage{bbm}
\usepackage{bm}
\usepackage{enumerate}
\usepackage[all]{xy}
\usepackage{natbib}
\usepackage{xcolor}
\usepackage{tikz}
\usepackage{braket}
\usepackage[shortlabels]{enumitem}
\newtheorem{lemma}{Lemma}
\newcommand{\im}{\mathrm{im}}
\title{Hodge Theory for Entanglement Cohomology}
\author{Christian Ferko, Eashan Iyer, Kasra Mossayebi, Gregor Sanfey}
\date{Source: arXiv:2410.12529v4 (CC BY 4.0)}
\begin{document}
\maketitle

\noindent\emph{Source and licence.} Hodge Theory for Entanglement Cohomology. Version arXiv:2410.12529v4 (CC BY 4.0), distributed by arXiv under the Creative Commons Attribution 4.0 International licence, http://creativecommons.org/licenses/by/4.0/, as recorded in the arXiv licence metadata for this exact version and shown on the abstract page https://arxiv.org/abs/2410.12529v4. Full licence notice: \texttt{licenses/arxiv-2410.12529-CC-BY-4.0.txt}.\par\smallskip

\noindent\textit{Editorial scope.} Counted unit: the article's lemma lemma\_harmonic\_rep (source0.tex lines 1208--1217). The article's complete original proof of it is delivered with it (source0.tex lines 1219--1269, one \texttt{proof} environment of 51 lines). No mathematics is written, repaired or re-derived: the statement and the proof are the article's own lines, in the article's own notation, and no retained line is substituted or re-flowed.  No further source-local context is needed: every cross-reference of every retained line resolves inside the two delivered spans.  Nothing was omitted from any retained span, so the package carries no omission record.  No retained line contains a citation, so the wrapper carries no bibliography and no bibliography entry.  No retained line contains a figure environment, an image-inclusion command or drawing code, so no image is delivered and the package declares no asset dependency.  Editorial formatting only, in addition: an AMS \texttt{amsart} preamble that restates the article's own declarations and macros used by the retained lines, namely the environments \texttt{lemma} on separate counters, together with the standard packages those lines need.  The retained environments print in this document as Lemma 1 in order of appearance, whereas the article prints the same items with the numbering of its own full document.  The complete native source of the article as distributed in the arXiv e-print for this exact version is bundled unchanged as \texttt{sources/166931/source0.tex}.
\begin{lemma}\label{lemma_harmonic_rep}
Consider a cochain complex
    %
\begin{align}
    \cdots \xrightarrow{d^{-1}} V_0 \xrightarrow{d^{0}} V_1 \xrightarrow{d^{1}} V_2 \xrightarrow{d^{2}} \ldots \, ,
\end{align}
%
in which each of the $V_i$ is a finite-dimensional Hilbert space with inner product $\langle \cdot \, , \, \cdot \rangle : V_i \times V_i \to \mathbb{C}$. Let $\delta^i : V_{i + 1} \to V_{i}$ be the $i$-th codifferential, which is the adjoint of the coboundary operator $d^i$ with respect to this inner product. Then in each cohomology class there is a unique element which is annihilated by $\delta$, and which is the representative with the smallest norm in its cohomology class. 

\end{lemma}

\begin{proof}
    One trivially has the orthogonal decomposition
    %
    \begin{align}\label{Vi_decomp}
        V_i = \mathrm{im} \left( d^{i-1} \right) \oplus \left( \mathrm{im} \left( d^{i-1} \right) \right)^\perp \, ,
    \end{align}
    %
    where $\perp$ denotes the orthogonal complement with respect to the inner product. But 
    %
    \begin{align}
        \left( \mathrm{im} \left( d^{i-1} \right) \right)^\perp = \ker \left( \delta^{i-1} \right) \, ,
    \end{align}
    %
    in finite-dimensional Hilbert spaces, since
    %
    \begin{align}
        \langle d^{i-1} \omega , \eta \rangle = \langle \omega, \delta^{i - 1} \eta \rangle \, ,
    \end{align}
    %
    and thus $\eta \in V_i$ is orthogonal to the image of $d^{i-1}$ if and only if it is annihilated by $\delta^{i - 1}$.

    Therefore the decomposition (\ref{Vi_decomp}) can be written as
    %
    \begin{align}\label{Vi_decomp_two}
        V_i = \im \left( d^{i-1} \right) \oplus \ker \left( \delta^{i - 1} \right) \, .
    \end{align}
    %
    Let $[\omega]$ be a cohomology class in $V_i$ with representative $\omega$, so that $d^i \omega = 0$. We decompose $\omega \in V_i$ according to (\ref{Vi_decomp_two}) as
    %
    \begin{align}
        \omega = d^{i-1} \lambda + \eta \, ,
    \end{align}
    %
    where $\lambda \in V_{i-1}$ and $\delta^{i-1} \eta = 0$. Note that, since
    %
    \begin{align}
        d^i \omega = d^i \left( d^{i-1} \lambda + \eta \right) = d^i \eta = 0 \, ,
    \end{align}
    %
    that $\eta$ is annihilated by both $d^i$ and $\delta^{i-1}$. This choice of $\eta$ is unique, by the uniqueness of orthogonal decompositions; said differently, any other $\eta' \neq \eta$ in the class $[\omega]$ will differ from $\eta$ by an exact form, which means that it would have a non-trivial projection onto $\im ( d^{i-1} )$, contradicting the orthogonal decomposition (\ref{Vi_decomp_two}).

    Finally, $\eta$ has the smallest norm in the cohomology class $[\omega]$ since for any $\lambda \in V_{i-1}$,
    %
    \begin{align}
        \langle \eta + d^{i-1} \lambda , \eta + d^{i-1} \lambda \rangle &= \langle \eta , \eta \rangle + 2 \langle \eta, d^{i-1} \lambda \rangle + \langle d^{i-1} \lambda , d^{i-1} \lambda \rangle \\
        &= \langle \eta , \eta \rangle + 2 \langle \delta^{i-1} \eta, \lambda \rangle + \langle d^{i-1} \lambda , d^{i-1} \lambda \rangle \nonumber \\
        &= \langle \eta , \eta \rangle + \langle d^{i-1} \lambda , d^{i-1} \lambda \rangle \nonumber \\
        &\geq \langle \eta , \eta \rangle  \, . \qedhere
    \end{align}
    %
\end{proof}

\end{document}
